\input{../tex-inputs/latex-leseansicht-vorspann}

\section[Gustav Schwarzkopf an Arthur Schnitzler, 14. 5. 1897]{L04284 Gustav Schwarzkopf an Arthur Schnitzler, 14. 5. 1897}
\nopagebreak\mylabel{L04284v}
\rehead{ }\normalsize\beginnumbering\briefempfaengerindex{Schnitzler, Arthur@\textsc{Schnitzler, Arthur}!zzzSchwarzkopf, Gustav@\emph{von Gustav Schwarzkopf}!1897-05-141@{14. 5. 1897}|(be}
\toendnotes[C]{\smallbreak\pagebreak[2]}
\correspDesc{Versand  durch Gustav Schwarzkopf am 14. 5. 1897 in Wien
\newline{}Übermittlung  am 14. 5. 1897 in Wien
\newline{}Erhalt  durch Arthur Schnitzler im Zeitraum [15. 5. 1897 – 19. 5. 1897?] in Paris}\toendnotes[C]{\smallbreak}
\Standort{CUL, Schnitzler, B 96a.}
\physDesc{Postkarte, 865 Zeichen
\newline{}Handschrift: schwarze Tinte, deutsche Kurrent
\newline{}Versand: Stempel: »\nobreak{}\oindex{I., Innere Stadt@\textbf{I., Innere Stadt}, \emph{Verwaltungsgebiet}|pwk}Wien 1/1, 14. 5. 97, 6–7N\nobreak{}«.  }\toendnotes[C]{\smallbreak}\pstart{}{\pb}\textsc{Mr. Arthur Schnitzler}\pend{}\pstart{}\textsc{Paris\oindex{Paris@\textbf{Paris}, \emph{Hauptstadt}|pw}}\pend{}\pstart{}\textsc{5, rue de Maubeuge\oindex{5, rue de Maubeuge@\textbf{5, rue de Maubeuge}, \emph{Wohngebäude}|pw}}\pend{}{\bigskip}\vspace{1em}
\pstart
           \raggedleft{}{\pb}14. V. 97\pend
           \vspace{0.5em}
\pstart
           Lieber Arthur! Ich ko{\geminationm}e nicht mehr
               zurück! So heißt es doch im Poker? 7 Seiten,{ }ſo viel wie 98 Zeilen, das iſt nicht zu
               überbieten, wenigſtens nicht von mir. Ich verſuche es auch gar nicht,{ }ſchreibe Ihnen
               nur, um Ihnen für Ihren{ }ſchönen \label{K_L04284-1v}\edtext{Brief}{\lemma{\textnormal{\emph{Brief}}}\Cendnote{\textnormal{XXXX Auszeichnungsfehler: Dokument L04111 nicht gefunden.}}}\label{K_L04284-1}{ }ſchön zu danken – we{\geminationn} ich ein Redacteur wäre, würde ich die beliebte
               »Indiskretion« begehen, ihn teilweiſe abzudrucken, »A. Sch. äußert{ }ſich in einem
               Brief aus Paris\oindex{Paris@\textbf{Paris}, \emph{Hauptstadt}|pw} über die neueſten Erſcheinungen«
               u. ſ. w.– und da{\geminationn} um Ihnen zu Ihrem \label{K_L04284-2v}\edtext{Geburtstag – er iſt
               doch morgen}{\lemma{\textnormal{\emph{Geburtstag … morgen}}}\Cendnote{\textnormal{Am 15. 5. 1897 wurde Schnitzler 35.}}}\label{K_L04284-2} oder an einem der nächſten Tage – zu gratuliren. Sie werden mit Recht
               vermuten, daß mich nur Bosheit und Schadenfreude dazu veranlaſſen. Ich ka{\geminationn} mir doch das Vergnügen nicht verſagen, es
               feſtzuſtellen, daß auch andere Leute älter werden. Viel Vergnügen in London\oindex{London@\textbf{London}, \emph{Hauptstadt}|pw} und auf Wiederſehen in Wien\oindex{Wien@\textbf{Wien}, \emph{Verwaltungsgebiet}|pw}\pend
           
\pstart
           Herzlichſt{\\[\baselineskip]}Ihr{\\[\baselineskip]}\spacefill\mbox{Gustav Schwz}\pend
           \leftskip=0em{}\selectlanguage{ngerman}\endnumbering\briefempfaengerindex{Schnitzler, Arthur@\textsc{Schnitzler, Arthur}!zzzSchwarzkopf, Gustav@\emph{von Gustav Schwarzkopf}!1897-05-141@{14. 5. 1897}|)be}\mylabel{L04284h}
\begin{anhang}
\end{anhang}\newcommand{\dateiname}{L04284}\newcommand{\titel}{Gustav Schwarzkopf an Arthur Schnitzler, 14. 5. 1897}\newcommand{\editorInnen}{Herausgegeben von Jahnke, SelmaMüller, Martin Anton}\input{../tex-inputs/latex-leseansicht-abspann}
